\section{Fuente Atorada en la Manguera}

\subsection*{Clasificación: I-A}

\subsubsection{Acordone una área de 25 metros a la redonda de donde se encuentra la fuente.}
Asegúrese de que el área esté correctamente señalizada y aislada para prevenir el acceso no autorizado.

\subsubsection{Consulte su dosímetro y póngalo en cero antes de iniciar las maniobras. Anote la dosis acumulada.}
Es fundamental llevar un registro preciso de la exposición a la radiación para garantizar la seguridad personal y cumplir con las normativas vigentes.

\subsubsection{Coloque el blindaje media caña adelante del contenedor sobre la manguera.}
\textbf{Atención:} \textit{Efectúe esta maniobra rápidamente}.
El uso de blindajes es crucial para protegerse de la radiación mientras se manejan fuentes radiactivas. Realice esta acción con celeridad para minimizar la exposición.

\subsubsection{Tomando la mayor distancia que sea posible, jale con el maneral al contenedor hasta que se observe en el detector que la fuente ha quedado blindada.}
Esta técnica ayuda a asegurar que la fuente radiactiva esté completamente cubierta y minimiza la exposición directa a la radiación.